\subsection{Sprint 0}

Para esta sprint, estavam previstas as configurações iniciais do repositório, o entendimento do projeto, a prototipação das telas, o conhecimento da equipe e a divisão dos squads. Da minha parte, fiquei responsável por dividir os squads junto aos Ages 4, compreender o termo de abertura do projeto e elencar perguntas. Além disso, busquei entender as fraquezas e forças dos integrantes do meu squad, pensando em como alocar os recursos da equipe da melhor maneira.

A equipe conseguiu entregar tudo o que estava previsto. O Figma foi majoritariamente explorado pelos Ages 1; o banco foi entregue pelos Ages 2; a arquitetura e a configuração ficaram a cargo dos Ages 3; e os épicos e as user stories foram elaborados pelos Ages 4. Como Ages 4, estive muito presente na gestão e na criação de tarefas no ClickUp, além de me responsabilizar pela apresentação dos épicos aos stakeholders.

Apesar das entregas realizadas, alguns integrantes da equipe não parecem tão engajados ou autônomos quanto o esperado. De momento, isso ainda não nos custou nada, mas é um problema que pode vir a complicar o trabalho no futuro.

Nem todo time precisa ser perfeito para realizar uma entrega: existiram partes da equipe que trabalharam muito bem de maneira autônoma e outras que apresentaram mais complicações. Cabe a um bom gestor tomar nota de problemas como esses e utilizar os recursos humanos de maneira eficiente.

Depois de entrevistar membros do meu squad, passei a ter as informações necessárias para entender como devo lidar com cada situação. Planejo procurar colegas em fases anteriores da Ages e oferecer ajuda, dado que, neste momento, o time tem mais a ganhar com as minhas dicas do que com a execução de uma task por conta própria. Planejo também manter uma conexão direta com os outros Ages 4, para conseguirmos operar bem como equipe.
